% GENERATED FILE. Edit canonical lessons, then run npm run generate:books.
% canonical-source-hash: fnv1a64:cc118b943c9459e7
% canonical-lessons: HI-C361-sodh, HI-C361-prayog, HI-C361-vaijnanik, HI-C361-mapna, HI-C361-prayogsala

\chapter{Research}
\label{ch:hi-research}
\hlchaptermodality{361}

\begin{chapteropening}
\textbf{By the end of this chapter:} \emph{I can say what a researcher is working on.}

Complete the last lesson of chapter 361: I can say what a researcher is working on.
\end{chapteropening}

\section[\texorpdfstring{\hi{शोध} (śodh)}{शोध (śodh)}]{\hi{शोध} (śodh) — research}

\label{lesson:HI-C361-sodh}



\begin{quote}
Before the new one: say the Hindi for to threaten, then the Hindi for to coax.
\end{quote}

\subsection*{You'll want to know: \hi{शोध}}
\textbf{\hi{शोध}} — \emph{śodh} — \textquotedblleft{}research\textquotedblright{}.

\begin{quote}
\hi{मीरा} \hi{पानी} \hi{पर} \hi{शोध} \hi{कर} \hi{रही} \hi{है।}
\end{quote}

\begin{quote}
\emph{Mīrā pānī par śodh kar rahī hai.}
\end{quote}

\begin{quote}
Mira is doing research on water.
\end{quote}

Masculine.

\begin{cousinweb}
\textbf{\hi{शोध}} (\emph{śodh}) is a tatsama, borrowed from Sanskrit \emph{śodha}, \textquotedblleft{}purification, cleansing\textquotedblright{} (Monier-Williams). Monier-Williams does not give the sense of research, which is usually taken as a modern Hindi one.
\end{cousinweb}

\begin{grammarlens}[title={What you've built}]
One more word for giving reasons and opinions.
\end{grammarlens}

\subsection*{Your turn}
\begin{itemize}
\raggedright
  \item \emph{Say it:} \emph{śodh}
  \item \emph{Say it:} \emph{śodh}, once more
  \item \emph{Say it:} \emph{Mīrā pānī par śodh kar rahī hai.} — Mira is doing research on water.
  \item \emph{Recall:} say the Hindi for to threaten, then the Hindi for to coax, then say \emph{śodh} again
\end{itemize}

\begin{tcolorbox}[breakable,colback=teal!4,colframe=teal!35!black,title={Before you move on}]
What does \hi{शोध} mean? (\textquotedblleft{}Research\textquotedblright{}.) Say it once more.
\end{tcolorbox}
\section[\texorpdfstring{\hi{प्रयोग} (prayog)}{प्रयोग (prayog)}]{\hi{प्रयोग} (prayog) — an experiment}

\label{lesson:HI-C361-prayog}



\begin{quote}
Before the new one: say the Hindi for to coax, then the Hindi for research.
\end{quote}

\subsection*{You'll want to know: \hi{प्रयोग}}
\textbf{\hi{प्रयोग}} — \emph{prayog} — \textquotedblleft{}an experiment\textquotedblright{}.

\begin{quote}
\hi{छात्रों} \hi{ने} \hi{कक्षा} \hi{में} \hi{एक} \hi{प्रयोग} \hi{किया।}
\end{quote}

\begin{quote}
\emph{chātroṁ ne kakṣā meṁ ek prayog kiyā.}
\end{quote}

\begin{quote}
The students did an experiment in class.
\end{quote}

Masculine.

\begin{cousinweb}
\textbf{\hi{प्रयोग}} (\emph{prayog}) is a tatsama, borrowed from Sanskrit \emph{prayoga}, \textquotedblleft{}use, application\textquotedblright{} (Monier-Williams).
\end{cousinweb}

\begin{grammarlens}[title={What you've built}]
One more word for giving reasons and opinions.
\end{grammarlens}

\subsection*{Your turn}
\begin{itemize}
\raggedright
  \item \emph{Say it:} \emph{prayog}
  \item \emph{Say it:} \emph{prayog}, once more
  \item \emph{Say it:} \emph{chātroṁ ne kakṣā meṁ ek prayog kiyā.} — The students did an experiment in class.
  \item \emph{Recall:} say the Hindi for to coax, then the Hindi for research, then say \emph{prayog} again
\end{itemize}

\begin{tcolorbox}[breakable,colback=teal!4,colframe=teal!35!black,title={Before you move on}]
What does \hi{प्रयोग} mean? (\textquotedblleft{}An experiment\textquotedblright{}.) Say it once more.
\end{tcolorbox}
\section[\texorpdfstring{\hi{वैज्ञानिक} (vaijñānik)}{वैज्ञानिक (vaijñānik)}]{\hi{वैज्ञानिक} (vaijñānik) — a scientist}

\label{lesson:HI-C361-vaijnanik}



\begin{quote}
Before the new one: say the Hindi for research, then the Hindi for an experiment.
\end{quote}

\subsection*{You'll want to know: \hi{वैज्ञानिक}}
\textbf{\hi{वैज्ञानिक}} — \emph{vaijñānik} — \textquotedblleft{}a scientist\textquotedblright{}.

\begin{quote}
\hi{वैज्ञानिक} \hi{नई} \hi{दवा} \hi{पर} \hi{काम} \hi{कर} \hi{रहे} \hi{हैं।}
\end{quote}

\begin{quote}
\emph{vaijñānik naī davā par kām kar rahe haiṁ.}
\end{quote}

\begin{quote}
Scientists are working on a new medicine.
\end{quote}

Masculine or feminine, by the person meant.

\begin{cousinweb}
\textbf{\hi{वैज्ञानिक}} (\emph{vaijñānik}) is a tatsama, borrowed from Sanskrit \emph{vaijñānika}, \textquotedblleft{}proficient, learned\textquotedblright{} (Monier-Williams). The sense of a scientist is a modern one that Monier-Williams does not give.
\end{cousinweb}

\begin{grammarlens}[title={What you've built}]
One more word for giving reasons and opinions.
\end{grammarlens}

\subsection*{Your turn}
\begin{itemize}
\raggedright
  \item \emph{Say it:} \emph{vaijñānik}
  \item \emph{Say it:} \emph{vaijñānik}, once more
  \item \emph{Say it:} \emph{vaijñānik naī davā par kām kar rahe haiṁ.} — Scientists are working on a new medicine.
  \item \emph{Recall:} say the Hindi for research, then the Hindi for an experiment, then say \emph{vaijñānik} again
\end{itemize}

\begin{tcolorbox}[breakable,colback=teal!4,colframe=teal!35!black,title={Before you move on}]
What does \hi{वैज्ञानिक} mean? (\textquotedblleft{}A scientist\textquotedblright{}.) Say it once more.
\end{tcolorbox}
\section[\texorpdfstring{\hi{मापना} (māpnā)}{मापना (māpnā)}]{\hi{मापना} (māpnā) — to gauge, to measure (precisely)}

\label{lesson:HI-C361-mapna}



\begin{quote}
Before the new one: say the Hindi for an experiment, then the Hindi for a scientist.
\end{quote}

\subsection*{You'll want to know: \hi{मापना}}
\textbf{\hi{मापना}} — \emph{māpnā} — \textquotedblleft{}to gauge, to measure (precisely)\textquotedblright{}.

\begin{quote}
\hi{दर्ज़ी} \hi{ने} \hi{मेरी} \hi{कमर} \hi{मापी।}
\end{quote}

\begin{quote}
\emph{darzī ne merī kamar māpī.}
\end{quote}

\begin{quote}
The tailor measured my waist.
\end{quote}

Past, with \hi{ने} (he or she …): \hi{उसने} \hi{कमर} \hi{मापी}, \emph{usne kamar māpī}.

\begin{cousinweb}
\textbf{\hi{मापना}} (\emph{māpnā}) is inherited from \emph{*māpyate}, \textquotedblleft{}is measured\textquotedblright{}, a form CDIAL reconstructs from its descendants (CDIAL 10054).
\end{cousinweb}

\begin{grammarlens}[title={What you've built}]
One more word for giving reasons and opinions.
\end{grammarlens}

\subsection*{Your turn}
\begin{itemize}
\raggedright
  \item \emph{Say it:} \emph{māpnā}
  \item \emph{Say it:} \emph{māpnā}, once more
  \item \emph{Say it:} \emph{darzī ne merī kamar māpī.} — The tailor measured my waist.
  \item \emph{Recall:} say the Hindi for an experiment, then the Hindi for a scientist, then say \emph{māpnā} again
\end{itemize}

\begin{tcolorbox}[breakable,colback=teal!4,colframe=teal!35!black,title={Before you move on}]
What does \hi{मापना} mean? (\textquotedblleft{}To gauge, to measure (precisely)\textquotedblright{}.) Say it once more.
\end{tcolorbox}
\section[\texorpdfstring{\hi{प्रयोगशाला} (prayogśālā)}{प्रयोगशाला (prayogśālā)}]{\hi{प्रयोगशाला} (prayogśālā) — a laboratory}

\label{lesson:HI-C361-prayogsala}



\begin{quote}
Before the new one: say the Hindi for a scientist, then the Hindi for to gauge.
\end{quote}

\subsection*{You'll want to know: \hi{प्रयोगशाला}}
\textbf{\hi{प्रयोगशाला}} — \emph{prayogśālā} — \textquotedblleft{}a laboratory\textquotedblright{}.

\begin{quote}
\hi{हमारे} \hi{स्कूल} \hi{में} \hi{एक} \hi{नई} \hi{प्रयोगशाला} \hi{है।}
\end{quote}

\begin{quote}
\emph{hamāre skūl meṁ ek naī prayogśālā hai.}
\end{quote}

\begin{quote}
There is a new laboratory in our school.
\end{quote}

Feminine.

\begin{cousinweb}
\textbf{\hi{प्रयोगशाला}} (\emph{prayogśālā}) is built from \emph{prayoga}, \textquotedblleft{}use, application\textquotedblright{}, and \emph{śālā}, \textquotedblleft{}a hall, a house\textquotedblright{} (Monier-Williams): a hall for experiments.
\end{cousinweb}

\begin{grammarlens}[title={What you've built}]
One more word for giving reasons and opinions.
\end{grammarlens}

\subsection*{Your turn}
\begin{itemize}
\raggedright
  \item \emph{Say it:} \emph{prayogśālā}
  \item \emph{Say it:} \emph{prayogśālā}, once more
  \item \emph{Say it:} \emph{hamāre skūl meṁ ek naī prayogśālā hai.} — There is a new laboratory in our school.
  \item \emph{Recall:} say the Hindi for a scientist, then the Hindi for to gauge, then say \emph{prayogśālā} again
\end{itemize}

\begin{tcolorbox}[breakable,colback=teal!4,colframe=teal!35!black,title={Before you move on}]
What does \hi{प्रयोगशाला} mean? (\textquotedblleft{}A laboratory\textquotedblright{}.) Say it once more.
\end{tcolorbox}
